\section{Benchmark Protocol and Environment}
\label{app:protocol}

\paragraph{Hardware and software.}
All measurements: one NVIDIA A100 80GB PCIe (dedicated, otherwise idle; driver
530.30) in a dual-socket host with two Intel Xeon Gold 6348 processors (Ice Lake,
28 cores and 56 threads each at 2.6\,GHz base, 112 logical CPUs, two NUMA nodes),
256\,GB of DDR4-3200 ECC memory on eight channels per socket, Linux 5.4 with the
\texttt{ondemand} frequency governor, GCC 9.4, PyTorch 1.13.1+cu117, CUDA 11.7,
Python 3.9.16. The three stages of the performance ladder (\cref{fig:ladder}) are
measured in one process under this stack: the pre-optimization iteration body is
preserved verbatim in the artifact (\texttt{opcount\_legacy.py}) and timed with the
same CUDA-event protocol as the compiled stages, and the same three stages re-measured
under PyTorch 2.5.1+cu121 are in \cref{tab:altexec}. CPU rows use 8
threads (\texttt{OMP\_NUM\_THREADS} and \texttt{MKL\_NUM\_THREADS} both set to~8)
pinned via
\texttt{taskset} to 8 whole physical cores on one NUMA node, with those cores' SMT
siblings excluded from the set (siblings share one core's execution units, so
counting them would overstate the arm's parallelism). LiteEFG 0.1.5 (single-threaded C++); OpenSpiel
Python CFR with exact best response. The multi-threaded CPU execution that the GPU
path is compared against is therefore this eight-thread compiled arm together with
the C++ library; \cref{fig:cpuscaling} reports how the compiled arm scales from one
to 28 threads.

\paragraph{Timing rules.}
GPU steady-state: CUDA events around a 1{,}000-iteration batch after 50
state-changing warmup iterations; the reported ms/iteration is the cost of iterations
51--1050. Training-call time (the ``Train'' columns of
\cref{tab:riverthreeway,tab:sg4threeway}): wall-clock of one training call with the
stated number of solver updates on a freshly constructed solver, so including the
call's three eager warmup iterations and graph capture,
but excluding game specification, compilation and solver construction (reported
separately in \cref{tab:overhead,tab:lifecycle}); a device
synchronize closes the training call before the evaluation bracket opens, so
asynchronous graph replays are never booked as evaluation time. The
warmup count differs because steady-state timing deliberately fills the captured
graph before measurement, whereas a training call measures the production setup
path. CPU and baseline
rows use \texttt{perf\_\allowbreak counter} wall-clock over the training call. Exploitability evaluation
time is excluded from all training times. The host is shared, so every reported number
comes from one serial gated campaign under the same protocol: 2{,}293 measurement cells
over 24.2 hours in two queues, a GPU queue pinned to one A100 and physical cores 48--55
(1{,}351 cells, 10.1 hours) and a CPU queue pinned to physical cores 0--7 (942 cells,
14.1 hours), covering the timing matrix, the ablations, the external baselines, the
convergence grids and the anytime curves.
Each cell is an independent process, with the 1-minute load average sampled before and
after it. A cell was not started above a load of 50 on the 112-core host; the realized
maximum was 7.4 in the GPU queue and 40.9 in the CPU queue, the latter from other
tenants on cores outside both pinned sets. Only two OpenSpiel cells ran above a load of
20, and the thread-scaling cells of \cref{fig:cpuscaling} additionally waited for a
load below 10. No cell was skipped or retried on that account.
Every speedup range and median quoted in the text is computed from the unrounded
cell medians by the table generators and enters the text as a generated macro, so
ratios of the rounded entries in \cref{tab:matrix} can differ from the quoted ones
in the last digit. Medians are over 20 repetitions for
GPU-CFR and the Ours CPU arm and 10 for the
vanilla-CFR control arm (regret matching without the $^{+}$ floor, linear averaging
retained) and for the baselines. Every archived result records the load at measurement time.

\paragraph{Operation and memory counters.}
The Aten operation counts of \cref{tab:suite} are collected with a
\texttt{TorchDispatchMode} interceptor, so they are load-immune and exactly
reproducible. \cref{tab:memory} reports peak GPU memory in binary megabytes from
the two \texttt{torch.cuda} counters, \texttt{max\_memory\_allocated} and
\texttt{max\_memory\_reserved}.

\paragraph{Solve protocol.}
CFR$^{+}$ (regret matching$^{+}$, linear averaging), simultaneous updates, float32,
seed 0. OpenSpiel runs 200 iterations of its vanilla CFR. Its three solvers
(CFR, CFR$^{+}$, DCFR) all fix \texttt{alternating\_updates=True}, so none of its
presets matches GPU-CFR's update order and we pick the one whose rule matches the
vanilla arm. Exploitability =
NashConv$/2$ for every backend. The update-rule study of \cref{tab:convergence} departs
from this protocol by design, sweeping four rules under both orders on the CPU at
$T{=}1{,}000$ in float32. It runs on the CPU because that device reproduces a solve
bitwise while the GPU does not, which is a requirement for ordering rules that sit close
together; \cref{app:correctness} quantifies the GPU spread. Every wall-clock number in
the paper, including the CPU arm of \cref{tab:matrix}, is unaffected. Quality comparisons across backend defaults use within-game wall-clock curves
(\cref{sec:quality}).

\begin{table*}[!htb]
\centering
\caption{Semantics of every timed arm. Order: simultaneous (S) or alternating (A,
one half-update per player per iteration). Averaging: uniform (U) or linear (L)
iterate weights. Every arm starts from uniform strategies and zero regrets; the
LiteEFG preset arm appears in \cref{app:extra}.}
\label{tab:semantics}
{
\setlength{\tabcolsep}{3pt}
\begin{tabular}{llccll}
\toprule
Arm & Rule & Order & Avg. & Precision & Execution \\
\midrule
GPU-CFR graph & CFR$^{+}$ & S & L & float32 & graph replay \\
GPU-CFR eager & CFR$^{+}$ & S & L & float32 & eager CUDA \\
GPU-CFR vanilla & CFR & S & L & float32 & graph replay \\
Ours CPU & CFR$^{+}$ & S & L & float32 & eager, 8 threads \\
Kim (2026) & CFR$^{+}$ & S & L & float32 & CuPy, eager \\
LiteEFG & CFR & S & U & float64 & C++, 1 thread \\
LiteEFG preset & CFR$^{+}$ & A & L & float64 & C++, 1 thread \\
OpenSpiel & CFR & A & U & float64 & Python, 200 it. \\
\bottomrule
\end{tabular}}
\end{table*}

\paragraph{Conversion validation.}
Public games are instantiated once through OpenSpiel and converted per backend.
Conversions are checked structurally (node, edge, and information-set counts match
across backends) and behaviorally: on the small games, exploitability curves and
their limiting values are consistent across backends, and Dark Hex 2$\times$2
reproduces its known first-player-win value ($v_0 = +1$) with exploitability
$\to 0$ under every backend.

\paragraph{Reproduction.}
Result files are JSONL with full environment echo (device, dtype, versions,
node/infoset counts, execution mode, load average). All tables and figures regenerate
from these files with the committed scripts. Campaign orchestration uses a separate
gated driver; the figure driver consumes only archived records and never launches a
solver. The 8-thread CPU arm was re-measured after the main campaign on whole physical
cores; both sets of records are archived, and the loader reads
the whole-core set. The artifact is provided as supplementary material: the solver,
the benchmark drivers, the archived JSONL records, the figure and table generators,
and the test suite, with a README that lists a smoke run of the test suite and the
full campaign commands together with the tolerances each check applies.
